\section{Parte 2: cómo hacer investigación en RL profundo}

\subsection{Actitud y metodología de investigación}

\begin{importantbox}{Recomendaciones de investigación de Chelsea Finn}
\begin{enumerate}
    \item \textbf{Partir del problema}: primero entender qué problema se quiere resolver y después pensar qué método usar. No «tener un martillo y buscar clavos».
    \item \textbf{Entender antes de innovar}: comprender a fondo las ventajas y desventajas de los métodos existentes; la innovación surge de forma natural.
    \item \textbf{Lo simple primero}: empezar por el método o el experimento más sencillo y aumentar la complejidad de manera gradual.
    \item \textbf{Construir intuición}: usar entornos sencillos para formar intuición sobre el comportamiento del algoritmo y luego generalizarla a escenarios complejos.
\end{enumerate}
\end{importantbox}

\subsection{Diseño experimental}

\begin{knowledgebox}{Buenas prácticas para experimentos de RL}
\begin{itemize}
    \item \textbf{Control de variables}: cambiar un solo factor a la vez.
    \item \textbf{Ejecuciones con varias semillas}: al menos de 3 a 5 semillas aleatorias.
    \item \textbf{Líneas base significativas}: elegir líneas base sólidas y no hombres de paja.
    \item \textbf{Experimentos de ablación}: verificar la contribución de cada decisión de diseño.
    \item \textbf{Herramientas de diagnóstico}: monitorear las curvas de aprendizaje, el error de estimación de la función de valor, la entropía de la política, etc.
\end{itemize}
\end{knowledgebox}

\begin{figure}[H]
\centering
\includegraphics[width=0.9\textwidth]{slides-images/slide-15.jpg}
\caption{El curso concreta «cómo escalar RL» en tres preguntas de diseño experimental: por lotes contra en línea, precisión de la función de valor y frecuencia de actualización. Estas preguntas determinan si un algoritmo puede escalar realmente a escenarios de horizonte largo y de modelos grandes.}
\end{figure}

\subsection{Depuración de algoritmos de RL}

Depurar algoritmos de RL es mucho más difícil que en el aprendizaje supervisado. Estrategias de depuración que propone la ponente:

\begin{enumerate}
    \item \textbf{Validar en entornos sencillos}: confirmar primero que el algoritmo funciona correctamente en CartPole/Pendulum.
    \item \textbf{Revisar la curva de recompensa}: la recompensa debe crecer de forma monótona (o aproximadamente monótona).
    \item \textbf{Revisar la función de valor}: la estimación de la función de valor debe acercarse al retorno real.
    \item \textbf{Visualizar el comportamiento de la política}: observar qué hace realmente el agente.
    \item \textbf{Revisar los gradientes}: asegurarse de que los gradientes no exploten ni se desvanezcan.
\end{enumerate}

\begin{warningbox}{Errores comunes al depurar RL}
\begin{itemize}
    \item Una scale inadecuada de la recompensa provoca que la función de valor diverja.
    \item Una elección inadecuada del factor de descuento $\gamma$.
    \item Omitir la normalización o estandarización de los datos.
    \item Una exploration insuficiente hace que la política converja a un óptimo local.
    \item La aleatoriedad puede ocultar un bug: aun con un bug, a veces RL logra aprender una política «aceptable».
\end{itemize}
\end{warningbox}

\subsection{Cómo elegir una línea de investigación}

La ponente ofrece recomendaciones para elegir una línea de investigación:
\begin{itemize}
    \item Enfocarse en \textbf{problemas realmente importantes} y no en temas de moda.
    \item Buscar la \textbf{brecha} entre la teoría y la práctica.
    \item Pensar entre disciplinas (NLP $\times$ RL, robotics $\times$ RL).
    \item Conversar con profesionales del área para conocer las necesidades reales.
\end{itemize}

\subsection{Resumen de la sección}

Una buena investigación en RL requiere una definición clara del problema, un diseño experimental riguroso y estrategias de depuración eficaces. Lo simple primero y el control de variables son la metodología central.
